\documentclass[12pt,a4paper,oneside]{report}
\usepackage[utf8]{inputenc}
\usepackage[T1]{fontenc}
\usepackage{lmodern}
\usepackage[margin=1in]{geometry}
\usepackage{setspace}
\usepackage{graphicx}
\usepackage{booktabs}
\usepackage{amsmath}
\usepackage[hidelinks]{hyperref}
\onehalfspacing
\newcommand{\pending}[1]{\par\noindent\textbf{Pending:} #1\par}
% Begin metadata
% Replace placeholders only with confirmed institutional/team details.
\newcommand{\thesistitle}{Input-Adaptive Teacher-to-Student Transition Timing for Efficient Diffusion Speech Synthesis}
\newcommand{\thesisauthors}{TBD: confirmed author names}
\newcommand{\thesisadviser}{TBD: adviser name}
\newcommand{\thesisinstitution}{TBD: institution}
\newcommand{\thesisprogram}{TBD: degree program}
\newcommand{\thesisdate}{October 2026}
% End metadata
\begin{document}
\begin{titlepage}
\centering
{\Large\thesisinstitution\par}
\vspace{1.8cm}
{\LARGE\bfseries\thesistitle\par}
\vspace{1.3cm}
{\large A working thesis manuscript\par}
\vspace{1cm}
{\large\thesisauthors\par}
\vfill
Adviser: \thesisadviser\par
Program: \thesisprogram\par
\vspace{0.7cm}
\thesisdate\par
\vspace{0.5cm}
\textbf{DRAFT --- research and institutional review pending}
\end{titlepage}
\pagenumbering{roman}
% Begin frontmatter/abstract
\chapter*{Abstract}
\addcontentsline{toc}{chapter}{Abstract}
This working manuscript investigates whether input-dependent teacher-guidance budgets can reduce the measured cost of diffusion speech synthesis while satisfying predefined speech-quality tolerances. The planned comparison includes fixed schedules, an SNR-only policy, and a hybrid policy whose extra computation must be accounted for. Existing repository telemetry provides development evidence and exposes reproducibility issues; it does not establish held-out effectiveness. The current workspace separates historical evidence, maintained analysis, and planned evaluation.
\pending{Rewrite the abstract after the protocol, held-out experiments, uncertainty analysis, and conclusions are complete. No final performance claim is made here.}
% End frontmatter/abstract
\tableofcontents
\listoftables
\listoffigures
\clearpage
\pagenumbering{arabic}
% Begin chapters/01-introduction
\chapter{Introduction}
\label{ch:introduction}
\section{Research problem}
The working research question is whether an inference-time policy can allocate a teacher-guidance budget according to input characteristics and reduce measured computational cost while satisfying predefined speech-quality and diversity tolerances on unseen inputs. This question is motivated by the internal source and telemetry audit~\cite{internal-audit-2026}.
\section{Objectives}
\begin{enumerate}
\item Reconcile the implementation's schedule labels, actual time coordinates, and forward-call accounting.
\item Reproduce a faithful baseline and compare tuned fixed schedules with a simple SNR-only candidate.
\item Determine whether a trajectory probe adds a meaningful held-out advantage after full cost accounting.
\item Evaluate intelligibility, speaker similarity, naturalness, failure cases, and across-seed diversity with speaker-level uncertainty.
\end{enumerate}
\section{Scope and limitations}
Duration prediction is held fixed and is outside the proposed contribution. Existing stored records are development evidence. Neither an abnormal teacher--student collision nor the necessity of a hybrid scheduler is established by that evidence. The proposed title and contribution remain subject to literature verification and adviser review.
\pending{Add institution-approved problem framing, stakeholder context, research questions, significance, and delimitations with verified sources.}
% End chapters/01-introduction
% Begin chapters/02-related-work
\chapter{Related Work}
\label{ch:related}
\section{Review plan}
The review will distinguish prior teacher-guided speech synthesis, input-dependent diffusion computation, speech-quality evaluation, and uncertainty estimation for repeated speaker measurements. For each source, the literature matrix will record the exact version, method, evidence, limitations, and relation to this thesis.
\section{Contribution boundary}
The internal audit identifies prior teacher-guided sampling and adaptive step allocation as boundaries that the novelty argument must address~\cite{internal-audit-2026}. A new threshold alone is not sufficient evidence of a novel contribution.
\pending{Read and verify primary publications before inserting bibliographic entries or specific prior-work claims. The internal audit citation is a provenance reference, not a substitute for a literature review.}
% End chapters/02-related-work
% Begin chapters/03-methodology
\chapter{Methodology}
\label{ch:method}
\section{Planned comparison}
The protocol will compare teacher-only and student-only references, a faithful guided baseline, development-tuned fixed schedules, an SNR-only policy, and a hybrid candidate. A retrospective oracle will be reported separately as offline headroom. Teacher-step counts, actual grid times, teacher and student forward calls, and end-to-end runtime will be recorded separately.
\section{Development replay}
The maintained replay uses the imported fixed-sweep records. Its SNR-only branch selects the historical three-teacher-step record for estimated SNR at least 18 dB and the four-teacher-step record otherwise. The historical labels are lookup identifiers rather than verified actual flow-time coordinates.

For this historical replay only, relative pass criteria are
\begin{align}
\mathrm{WER}_{\mathrm{out}} &\leq \max(12.5,\mathrm{WER}_{\mathrm{base}}+1),\\
\mathrm{SIM}_{\mathrm{out}} &\geq \mathrm{SIM}_{\mathrm{base}}-0.02.
\end{align}
WER is expressed in percentage points. The strict variant fixes the WER ceiling at 12.5. These operational rules are not universal quality standards and do not establish strict non-degradation~\cite{internal-audit-2026}.
\section{Held-out evaluation plan}
Speakers will be separated before tuning, and thresholds will be frozen on development data. Matched noise and latent seeds will support paired comparisons. The protocol will specify meaningful quality margins and speaker-level variability before choosing test size. Timing will include warmup, device synchronization where relevant, and all scheduler probe or restart work.
\pending{Freeze data manifests, evaluator definitions, quality margins, sample size, statistical procedure, listening protocol, and implementation revision before collecting final results.}
% End chapters/03-methodology
% Begin chapters/04-results
\chapter{Results and Discussion}
\label{ch:results}
\section{Available development evidence}
The October 5 audit reports that the stored timing records do not establish a hybrid advantage over the simpler SNR-only replay~\cite{internal-audit-2026}. The new workspace contains an executable replay and hashed inputs under experiment EXP-000. Its output is descriptive arithmetic on stored records, not an independent runtime measurement.

\section{Planned reporting}
Report per-policy quality and compute estimates with uncertainty, absolute failures, baseline-relative changes, and speaker-cluster analysis. Include negative outcomes and inspect catastrophic ASR or synthesis cases. Distinguish scheduler overhead from generation cost and cross-seed diversity from single-waveform pitch variation.
\pending{Insert independently validated held-out results, tables, figures, confidence intervals, and listening outcomes. Link each quantitative claim to a completed experiment record.}
% End chapters/04-results
% Begin chapters/05-conclusion
\chapter{Conclusion and Future Work}
The workspace currently supports source provenance, reproducible development-record analysis, manuscript drafting, and evidence-based progress records. The research question remains open. Final conclusions will depend on baseline reconciliation and a frozen held-out evaluation.
\pending{State only conclusions supported by completed experiments. Explain whether adaptivity and the hybrid probe added meaningful value, including a null or negative outcome if that is what the evidence shows.}
% End chapters/05-conclusion
\begin{thebibliography}{1}
\bibitem{internal-audit-2026}
\emph{Thesis Direction Audit: 5 October 2026}. October 2026.
Internal repository audit. Archived at
\texttt{research/audits/2026-10-05/FINDINGS.md}; source and telemetry analysis,
not independent synthesis evaluation.
\end{thebibliography}
\appendix
% Begin appendices/reproducibility
\chapter{Reproducibility Record}
The repository source manifest identifies imported file paths and SHA-256 hashes. Experiment EXP-000 records the timing replay command, evaluator configuration, inputs, and limitations. Future model experiments must additionally record checkpoint hashes, manifests, source revisions, environment exports, hardware, seeds, complete cost accounting, and output artifact locations.

The manuscript uses a provisional report layout. Institutional formatting, author information, adviser details, and submission requirements remain unconfirmed.
% End appendices/reproducibility
\end{document}
